\section{Crossing order and Hall--skein normal forms}
\label{sec:crossing}

We now compare the exact Hall product with polygon smoothing.  The key point is
not that the Hall coefficients equal skein coefficients before quotienting;
they generally do not.  The key point is triangularity: the split crossing
object is the unique top term, and every non-split middle term lies lower in a
well-founded smoothing order.

\subsection{Extension middle terms in type $A$}

Let $M_\gamma,M_\delta$ be the reduced morphism objects corresponding to
polygon diagonals $\gamma,\delta$.  An exact extension in $\MM$
\[
 0\longrightarrow M_\delta\longrightarrow E\longrightarrow M_\gamma
 \longrightarrow0
\]
induces a triangle in the two-term homotopy category and, after applying the
polygon realization, a triangle in the cluster category.  The two directions
of the $2$-Calabi--Yau extension space correspond to the two smoothings of a
crossing.

\begin{theorem}[Type-$A$ extension--smoothing triangularity]
\label{thm:extension-smoothing}
Let $X,Y$ be reduced objects of $\MM$ and let $E$ be the middle term of an
exact extension of $X$ by $Y$.
\begin{enumerate}[label=(\roman*),leftmargin=2.3em]
\item If the extension is split, then $E\simeq X\oplus Y$.
\item If the extension is non-split, the reduced polygon diagram of $E$ is
      obtained from the superposition $D(X)\cup D(Y)$ by smoothing at least
      one crossing, followed by removal of contractible summands and polygon
      isotopy.
\item Consequently every non-split middle term has strictly fewer crossings
      than the inherited superposition diagram of $X\oplus Y$.
\end{enumerate}
For indecomposable arcs crossing once, the two possible reduced middle terms
in the two extension directions are precisely
$\gamma\#_0\delta$ and $\gamma\#_\infty\delta$, up to contractible summands
and boundary coefficient factors.
\end{theorem}

\begin{proof}
The exact-to-homotopy identification
\[
 \Ext^1_{\MM}(X,Y)
 \simeq\Hom_{K^b(\proj A)}(X,\Sigma Y)
\]
turns an exact extension into its mapping-cone triangle.  Under the type-$A$
polygon model, indecomposable two-term objects are arcs and nonzero morphisms
of degree one are represented by crossings.  Canakci--Schroll describe bases
of extension spaces for the gentle Jacobian algebras of unpunctured surfaces
and compute the middle terms by smoothing the corresponding crossings
\cite{CanakciSchroll}.  In a polygon there are no bands and no puncture
phenomena, so every indecomposable extension is of this form.

For direct sums, write an extension class as a matrix of its indecomposable
components.  Successive pullback--pushout decompositions express its middle
term as an iterated extension along the nonzero matrix entries.  Applying the
indecomposable smoothing description at each nonzero entry gives the claimed
simultaneous smoothing picture.  A local smoothing removes the chosen
geometric crossing.  Keeping the inherited embedding, it introduces no new
crossing; passing to the unique minimal isotopy classes of arcs in a disk can
only remove further crossings.  The extension is non-split exactly when at
least one matrix entry is nonzero, so the decrease is strict.
\end{proof}

\begin{remark}
The theorem is special to the surface/gentle finite-type setting.  For a
general Jacobian algebra, extension middle terms need not admit a two-smoothing
model.  This is the precise reason the global theorem below is restricted to
type $A$.
\end{remark}

\subsection{A well-founded order}

For a reduced object $M$, let $r(M)$ be the number of indecomposable summands,
counted with multiplicity, and let $\crs(M)$ be the crossing number from
\eqref{eq:crossing-number}.  Fix also a total refinement of the degeneration
order on each finite representation variety.  Define
\[
 E\prec M
\]
if, in lexicographic order,
\begin{enumerate}[label=(\alph*),leftmargin=2.2em]
\item $E$ has smaller crossing number than $M$; or
\item the crossing numbers agree and $E$ is strictly more generic in the
      extension-degeneration order; or
\item both agree and the fixed total refinement places $E$ first.
\end{enumerate}
Within a fixed term class and bounded number of summands this is a
well-founded order.  The only property used later is that every non-split
middle term in \cref{thm:extension-smoothing} lies below the split middle term.

For Hall basis objects $X,Y$, write the cocycle-twisted product as
\begin{equation}
 \eps_X\star_S\eps_Y
 =h^{\rm sp}_{X,Y}\eps_{X\oplus Y}
  +\sum_{E\prec X\oplus Y}h^E_{X,Y}\eps_E.
 \label{eq:Hall-triangular-product}
\end{equation}
The split coefficient $h^{\rm sp}_{X,Y}$ is nonzero.  We work over
$\mathbb Q(\omega)$, so it is invertible.

\begin{corollary}[Crossing-leading Hall products]
\label{cor:crossing-leading}
If $\gamma$ and $\delta$ cross, then
\begin{equation}
 \eps_\gamma\star_S\eps_\delta
 =h^{\rm sp}_{\gamma,\delta}\eps_{\gamma\oplus\delta}
  +\sum_{E\prec\gamma\oplus\delta}
      h^E_{\gamma,\delta}\eps_E,
 \qquad h^{\rm sp}_{\gamma,\delta}\ne0,
 \label{eq:crossing-Hall-expansion}
\end{equation}
and every $E$ in the second sum has smaller crossing number.
\end{corollary}

\subsection{Hall--skein defects and oriented reductions}

For an ordered crossing pair define
\begin{equation}
 \mathcal R_{\gamma,\delta}
 :=\eps_\gamma\star_S\eps_\delta
   -\omega^{a(\gamma,\delta)}
      \eps_{\gamma\#_0\delta}
   -\omega^{b(\gamma,\delta)}
      \eps_{\gamma\#_\infty\delta}.
 \label{eq:Hall-skein-defect}
\end{equation}
The symbols on the right denote the normalized direct-sum classes of the two
smoothing components, with boundary components interpreted as frozen units.
Let
\begin{equation}
 \Jsk=
 \left\langle
   \mathcal R_{\gamma,\delta}:
   \gamma,\delta\text{ an ordered crossing pair}
 \right\rangle_{\rm two\text{-}sided}
 \subseteq\Hred.
 \label{eq:skein-ideal}
\end{equation}

Substituting \eqref{eq:crossing-Hall-expansion} into
\eqref{eq:Hall-skein-defect} gives an oriented rule
\begin{equation}
\begin{aligned}
 \eps_{\gamma\oplus\delta}
 \longmapsto
 (h^{\rm sp}_{\gamma,\delta})^{-1}
 \Big(&\omega^{a(\gamma,\delta)}
       \eps_{\gamma\#_0\delta}
      +\omega^{b(\gamma,\delta)}
       \eps_{\gamma\#_\infty\delta}\\
     &-\sum_{E\prec\gamma\oplus\delta}
       h^E_{\gamma,\delta}\eps_E\Big).
\end{aligned}
\label{eq:basic-Hall-straightening}
\end{equation}
Every term on the right has smaller crossing number.  Multiplying the defect
on the left and right by Hall classes and expanding triangularly yields the
same kind of rule for any object containing $\gamma\oplus\delta$ as a direct
summand.  The leading coefficient remains invertible.

\begin{lemma}[Termination]
\label{lem:termination}
Every sequence of oriented Hall--skein reductions terminates after finitely
many steps in a linear combination of non-crossing Hall basis objects.
\end{lemma}

\begin{proof}
Each reduction smooths a chosen geometric crossing in the inherited polygon
diagram.  By \cref{thm:extension-smoothing}, all Hall correction terms also
have strictly fewer crossings than the leading object.  Hence the
nonnegative integer $\crs$ decreases at every step.
\end{proof}

\subsection{The filtered Diamond lemma}

Termination alone does not guarantee a basis: different choices of crossing
could a priori give different non-crossing expressions.  We now prove
confluence.

\begin{theorem}[Hall--skein Diamond lemma]
\label{thm:Hall-skein-diamond}
The reduction system \eqref{eq:basic-Hall-straightening}, including all its
left and right Hall contexts, is confluent.  Consequently every class in
$\Hred/\Jsk$ has a unique normal form which is a linear combination of the
classes
\[
 \eps_D,
 \qquad D\in\mathfrak N(P).
\]
These non-crossing classes are linearly independent in the quotient.
\end{theorem}

\begin{proof}
We apply the filtered version of Bergman's Diamond lemma with the
well-founded crossing order.  It is enough to resolve critical ambiguities.
There are two types.

\smallskip
\noindent\emph{Disjoint crossings.}
If two chosen crossing disks are disjoint, the local smoothings commute.  On
the Hall side the two corresponding substitutions are made in disjoint
factors of an associative product.  Their difference is a sum of terms in
which both chosen crossings have been removed, so the two reductions agree
after lower reductions.

\smallskip
\noindent\emph{Overlapping crossings.}
Suppose the two reductions share an arc or occur in an iterated product of
three Hall generators.  Both reduction paths start from the same associative
Hall word.  Replace one crossing product by the right-hand side of
\eqref{eq:Hall-skein-defect}, and do the same in the other order.  The
resulting difference belongs to $\Jsk$, while cancellation of the common top
Hall word leaves only terms of smaller crossing number.  By induction on the
crossing number, every such lower element of $\Jsk$ reduces to zero.  At the
geometric level, this induction is exactly the local confluence of the polygon
skein relation from \cref{thm:skein-basis}; endpoint-state ambiguities are
resolved by the boundary exchange relations already imposed in the central
coefficient system.

The base case has crossing number zero.  No leading term of a Hall--skein
defect is non-crossing, so there is no reduction and no nonzero ambiguity in
that degree.  Thus all critical pairs are resolvable.  The Diamond lemma gives
unique normal forms and identifies the irreducible Hall basis classes with the
non-crossing multidiagonals.  In particular, no nontrivial linear combination
of irreducible classes belongs to $\Jsk$.
\end{proof}

\begin{remark}[Where Hall associativity is used]
The skein relation alone proves confluence in the free diagram algebra.  The
additional Hall correction terms in
\eqref{eq:basic-Hall-straightening} could have destroyed it.  Associativity of
$\star_S$ forces the difference between two Hall-corrected overlap reductions
to lie in the same two-sided ideal after cancellation of the common top word;
strict crossing decrease then closes the induction.  This is the essential
new algebraic step.
\end{remark}

\subsection{A basis theorem for the quotient}

\begin{corollary}[Rigid normal basis]
\label{cor:rigid-normal-basis}
The classes
\begin{equation}
 \left\{\overline\eps_D:D\in\mathfrak N(P)\right\}
 \label{eq:quotient-rigid-basis}
\end{equation}
form a $\mathbb Q(\omega)$-basis of $\Hred/\Jsk$.  Under the polygon
category correspondence, this basis is indexed by all rigid direct sums,
including arbitrary multiplicities of compatible indecomposables.
\end{corollary}

\begin{proof}
Spanning is \cref{lem:termination}; independence and uniqueness are
\cref{thm:Hall-skein-diamond}.  A multidiagonal is non-crossing exactly when
the corresponding direct sum has no self-extension.
\end{proof}
